% W4i26_P_updates.tex
% DFD 17-Agent Research Programme — Iteration 26
% Agents 6 (P1-P4 Dead Patents), 7 (P5-P7 Alive Patents), 8 (P8-P11 SQMS/SRF)
% Date: 2026-03-21

\documentclass[11pt,a4paper]{article}
\usepackage[utf8]{inputenc}
\usepackage{amsmath,amssymb,amsfonts}
\usepackage{hyperref}
\usepackage{booktabs}
\usepackage{longtable}
\usepackage{geometry}
\geometry{margin=2.5cm}

\title{DFD i26 Patent \& Technology Updates\\Agents 6, 7, 8}
\author{DFD 17-Agent Research Programme}
\date{2026-03-21}

\begin{document}
\maketitle
\tableofcontents
\newpage

%=============================================================================
\section{Agent 6: Dead Patents P1--P4 (Metamaterial/Cloaking)}
%=============================================================================

\subsection{Mandate}
Monitor metamaterial and cloaking technology developments relevant to dead patents P1--P4. Identify any revival opportunities from new technology breakthroughs.

\subsection{Patent Status Recap}
\begin{itemize}
\item \textbf{P1--P4}: All DEAD. Metamaterial cloaking, gradient-index optics, refractive shielding concepts.
\item Core limitation: practical metamaterial manufacturing remains cost-prohibitive for the scales envisioned in these patents.
\end{itemize}

\subsection{New Literature}

\begin{enumerate}

\item \textbf{Ren et al.\ (2025), Sci.\ Rep.\ 15, 11501} --- ``Large-scale underwater acoustic cloak based on honeycomb lattice pentamode metamaterials.'' Demonstrates acoustic cloaking at 1--3 kHz with total scattering cross-section reduced by factor 2. Honeycomb pentamode structure enables scalability. \textbf{Grade: C.} Large-scale acoustic cloaking demonstrates progress but remains narrowband. Not directly applicable to DFD's refractive-field concept, which operates on gravitational rather than acoustic waves.

\item \textbf{Xia et al.\ (2025), PNAS 122, e2502036122} --- ``Combinatorial asymmetric acoustic metamaterials with real-time programmability.'' Ultrascalable framework for programmable acoustic metamaterials. Real-time reconfiguration of wave-manipulation properties. \textbf{Grade: C.} Programmable metamaterials could in principle enable tunable refractive-index profiles, but the gap between acoustic and gravitational applications remains vast.

\item \textbf{Singh et al.\ (2025), Nat.\ Commun.\ Eng.\ s44172-025-00470-x} --- ``Bio-inspired acoustic metamaterials for traffic noise control: bridging the gap with machine learning.'' ML-optimized metamaterial design. \textbf{Grade: C.} Demonstrates AI-driven design methodology that could accelerate exploration of gradient-index structures, but no gravitational application.

\item \textbf{Iftikhar et al.\ (2024), ACS Appl.\ Mater.\ Interfaces 4c04486} --- ``AI-Based Metamaterial Design.'' Comprehensive review of machine learning methods for metamaterial inverse design. Generative models (GANs, VAEs) for novel structures. \textbf{Grade: C.} The AI design methodology is maturing rapidly and could in principle be applied to gradient-index gravitational lens design, but P1--P4 remain dead from a practical standpoint.

\end{enumerate}

\subsection{Revival Assessment}
\textbf{P1--P4 remain DEAD.} No breakthrough in 2025--2026 bridges the gap between acoustic/EM metamaterials and the gravitational regime. The refractive-index analogy ($n = e^\psi$) is intellectually connected but practically unrelated to current metamaterial technology.

\textbf{Long-term watch}: AI-driven metamaterial design (Iftikhar et al.) could accelerate discovery of novel gradient-index structures if materials science advances enable gravitational-frequency operation.

%=============================================================================
\section{Agent 7: Alive Patents P5--P7 (GW Detection)}
%=============================================================================

\subsection{Mandate}
Track Einstein Telescope (ET), Cosmic Explorer (CE), and LISA Pathfinder legacy. Assess implications for DFD's alive gravitational-wave patents P5, P7.

\subsection{Patent Status Recap}
\begin{itemize}
\item \textbf{P5}: ALIVE. Refractive-field GW detector enhancement concept.
\item \textbf{P7}: ALIVE. Scalar-field perturbation measurement via GW interferometry.
\item Core value: DFD predicts $\alpha_T = 0$ (GW speed = $c$) but with scalar-mode signatures detectable by next-gen interferometers.
\end{itemize}

\subsection{New Literature}

\begin{enumerate}

\item \textbf{Branchesi et al.\ (2025), arXiv:2505.11033} --- ``Einstein Telescope and Cosmic Explorer.'' Comprehensive review of 3G GW detector science case. ET site decision expected 2026 (Sardinia vs.\ Euregio Meuse-Rhine). Construction 2028, operation 2035. CE: 40 km arms, NSF-endorsed 2024, operational mid-2030s. Sensitivity: factor 8 improvement over LIGO/Virgo/KAGRA, extended to $<$10 Hz. \textbf{Grade: A.} P5 and P7 patent value increases as 3G detectors approach. The sub-10 Hz band is where DFD scalar breathing modes are predicted to be detectable.

\item \textbf{ET Governance (2025--2026)} --- Netherlands, Belgium, Germany, Italy all declare ET a national priority. German coalition agreement highlights ET. Site announcement expected 2026. \textbf{Grade: B.} Political support reduces timeline risk for P5/P7 exploitation window.

\item \textbf{de Rham \& Tolley (2025), arXiv:2511.19762} --- ``The Speed of Gravity and the Fate of Dark Energy.'' Reviews GW170817 constraints on dark energy models. Disformal couplings constrained: $M_\gamma \gtrsim 10$ MeV. Eliminates quartic galileon. \textbf{Grade: B.} DFD automatically satisfies this constraint ($\alpha_T = 0$), which eliminates many competitor theories. P7's scalar-mode concept becomes more distinctive as tensor-speed alternatives are ruled out.

\item \textbf{Takahashi et al.\ (2025), arXiv:2509.05027} --- ``Unique GW signatures of GLPV scalar-tensor theories.'' Induced GW spectrum $\propto f^5$ for GLPV models. New scalar-scalar-tensor interaction unique to beyond-Horndeski. \textbf{Grade: B.} DFD's disformal sector is related to GLPV. If DFD induces a characteristic $f^5$ spectral signature, this would be a smoking gun for P7's scalar-mode detection concept.

\end{enumerate}

\subsection{P5/P7 Exploitation Timeline}
\begin{center}
\begin{tabular}{lll}
\toprule
\textbf{Year} & \textbf{Milestone} & \textbf{Patent Relevance} \\
\midrule
2026 & ET site decision & P5/P7 licensing discussions \\
2028 & ET construction begins & P5 detector enhancement IP \\
2032 & LiteBIRD launch & Complementary $\beta=0$ test \\
2035 & ET/CE operational & P7 scalar-mode detection \\
\bottomrule
\end{tabular}
\end{center}

\textbf{Grade: A (P5/P7 value increasing).}

%=============================================================================
\section{Agent 8: P8--P11 (SQMS, SRF, Quantum Sensing)}
%=============================================================================

\subsection{Mandate}
Track SQMS/SRF developments, quantum sensing with cavities, transmon-qubit $\psi$-probe concept. P9 (alive), P11 (alive).

\subsection{Patent Status Recap}
\begin{itemize}
\item \textbf{P8}: DEAD. General SRF cavity concept.
\item \textbf{P9}: ALIVE. DFD-specific Q-ratio prediction ($Q_\psi / Q_\mathrm{BCS} = 0.52 \pm 0.08$ vs.\ BCS $3.60 \pm 0.10$, $>$30$\sigma$ discrepancy).
\item \textbf{P10}: DEAD. Generic quantum sensing.
\item \textbf{P11}: ALIVE. Transmon-qubit $\psi$-probe for local scalar field measurement.
\end{itemize}

\subsection{New Literature}

\begin{enumerate}

\item \textbf{Kalia, Liu et al.\ (2026), Phys.\ Rev.\ Lett.\ 136, 111802; arXiv:2510.02427} --- ``Improved Dark Photon Sensitivity from the Dark SRF Experiment.'' Refined exclusion bounds from Dark SRF pathfinder run at Fermilab SQMS. Order of magnitude improvement via better frequency-instability modeling. World-leading non-DM dark photon constraint below 6 $\mu$eV. Best lab-based photon mass limit: $m_\gamma < 2.9\times10^{-48}$ g. \textbf{Grade: A.} Directly validates SRF cavity sensitivity at levels relevant to P9's Q-ratio measurement. The frequency-instability modeling improvements are transferable to DFD's proposed $\psi$-field sensitivity analysis.

\item \textbf{Dark SRF Phase II plans (2025--2026)} --- SQMS planning next-generation run with improved cavity quality factors ($Q > 10^{11}$) and quantum-limited amplifiers. Target: wavelike DM dark photon sensitivity at $\epsilon < 10^{-16}$. \textbf{Grade: B.} The cavity infrastructure being developed for dark photon searches could be repurposed for P9/P11 $\psi$-field measurements with minimal modifications.

\item \textbf{SRF2025 Conference (Sept 2025), Proceedings} --- 22nd International Conference on RF Superconductivity. Reports on nitrogen-doped Nb cavities achieving $Q_0 > 5\times10^{10}$ at 2 K. New Nb$_3$Sn cavities with $Q_0 > 10^{10}$ at 4.2 K. \textbf{Grade: B.} Higher $Q$ cavities narrow the gap to the sensitivity needed for DFD's predicted anomalous Q-ratio.

\item \textbf{SQMS Quantum Networking (2025), ScienceDaily Oct 29} --- New quantum sensor network concept linking SRF cavities for correlated dark matter searches. Distributed sensing over km baselines. \textbf{Grade: B.} A networked SRF array could implement DFD's gradient measurement concept: measuring $\nabla\psi$ between spatially separated cavities to detect the scalar field gradient.

\end{enumerate}

\subsection{P9 Q-Ratio Update}

DFD prediction: $Q_\psi/Q_\mathrm{BCS} = 0.52 \pm 0.08$. This requires measuring the intrinsic quality factor of a high-purity Nb cavity with precision $\delta Q/Q < 10^{-3}$ while controlling for:
\begin{itemize}
\item Magnetic flux trapping ($\to$ nitrogen doping reduces this)
\item Two-level system losses ($\to$ surface treatment protocols from SRF2025)
\item Multipactor and field emission ($\to$ Dark SRF protocols)
\end{itemize}

The Dark SRF team's frequency-instability modeling (Kalia et al.\ 2026) directly addresses the dominant systematic for Q-ratio measurements. \textbf{P9 is experimentally closer than at any prior iteration.}

\subsection{P11 Transmon-Qubit $\psi$-Probe}

The transmon qubit coupled to an SRF cavity measures the effective cavity frequency to $\sim$Hz precision. DFD predicts a $\psi$-field-dependent frequency shift:
\begin{equation}
\delta\omega / \omega = -\psi \approx -\Phi_N / c^2 \approx -7\times10^{-10} \text{ (Earth surface)}.
\end{equation}
This corresponds to $\delta\omega \sim 4$ Hz at $\omega/(2\pi) \sim 6$ GHz. The SQMS quantum networking concept (distributed SRF cavities) could measure the \emph{gradient} $\nabla\psi$ between two locations, providing a differential measurement that cancels common-mode noise.

\textbf{P11 Grade: B+ (upgraded). Quantum network concept provides new measurement strategy.}

%=============================================================================
\section{Patent Portfolio Summary}
%=============================================================================

\begin{center}
\begin{tabular}{lcll}
\toprule
\textbf{Patent} & \textbf{Status} & \textbf{i26 Development} & \textbf{Grade} \\
\midrule
P1--P4 & DEAD & AI metamaterial design advancing, but gap vast & C \\
P5 & ALIVE & ET site 2026, construction 2028 & A \\
P7 & ALIVE & GLPV $f^5$ signature, GW speed constraints & A \\
P8 & DEAD & No change & --- \\
P9 & ALIVE & Dark SRF sensitivity improved $10\times$ & A \\
P10 & DEAD & No change & --- \\
P11 & ALIVE & Quantum network concept for $\nabla\psi$ & B+ \\
\bottomrule
\end{tabular}
\end{center}

%=============================================================================
\section{New Literature Master Table (Agents 6, 7, 8)}
%=============================================================================

\begin{longtable}{p{3cm}p{3cm}p{5cm}cc}
\toprule
\textbf{Reference} & \textbf{arXiv/DOI} & \textbf{Key Result} & \textbf{Agent} & \textbf{Grade} \\
\midrule
\endhead
Ren+ 2025 & Sci.\ Rep.\ 15, 11501 & Large-scale acoustic cloak, pentamode & 6 & C \\
Xia+ 2025 & PNAS 122, e2502036122 & Programmable acoustic metamaterials & 6 & C \\
Singh+ 2025 & Nat.\ Commun.\ Eng.\ & ML acoustic metamaterial design & 6 & C \\
Iftikhar+ 2024 & ACS AMI 4c04486 & AI-based metamaterial design review & 6 & C \\
Branchesi+ 2025 & 2505.11033 & ET/CE comprehensive science case & 7 & A \\
de Rham+ 2025 & 2511.19762 & GW speed constraints on disformal & 7 & B \\
Takahashi+ 2025 & 2509.05027 & GLPV $f^5$ GW signature & 7 & B \\
Kalia+ 2026 & PRL 136, 111802 & Dark SRF $10\times$ improved, best photon mass & 8 & A \\
SRF2025 Proceedings & Conf.\ Proc.\ & $Q > 5\times10^{10}$ Nb cavities & 8 & B \\
SQMS Network 2025 & ScienceDaily & Distributed SRF quantum sensing & 8 & B \\
\bottomrule
\end{longtable}

\end{document}
